\documentclass[12pt]{article}

\usepackage[T1]{fontenc}
\usepackage{lmodern}
\usepackage{microtype}
\usepackage[left=1in,right=1in,top=0.82in,bottom=0.84in]{geometry}
\usepackage{setspace}
\usepackage{amsmath,amssymb,mathtools}
\usepackage{mathrsfs}
\usepackage{fancyhdr}
\usepackage[hidelinks]{hyperref}

\hypersetup{
  pdftitle={How the Proof Works},
  pdfauthor={Thomas Birnie}
}

\setstretch{1.18}
\setlength{\parindent}{0pt}
\setlength{\parskip}{0.62\baselineskip}
\setlength{\abovedisplayskip}{13pt plus 3pt minus 2pt}
\setlength{\belowdisplayskip}{13pt plus 3pt minus 2pt}
\setlength{\abovedisplayshortskip}{11pt plus 2pt minus 2pt}
\setlength{\belowdisplayshortskip}{11pt plus 2pt minus 2pt}
\setlength{\jot}{9pt}
\setlength{\emergencystretch}{3em}
\allowdisplaybreaks[2]
\raggedright
\raggedbottom
\clubpenalty=10000
\widowpenalty=10000
\displaywidowpenalty=10000

\pagestyle{fancy}
\fancyhf{}
\renewcommand{\headrulewidth}{0pt}
\fancyfoot[C]{\footnotesize\thepage}

\newcommand{\R}{\mathbb R}
\newcommand{\I}{\mathcal I}
\newcommand{\cR}{\mathcal R}
\newcommand{\Szero}{\mathscr S_0}

\begin{document}

{\Large\bfseries How the Proof Works\par}
\vspace{0.16\baselineskip}
{\small Thomas Birnie\par}

\vspace{0.55\baselineskip}

For the global three-dimensional Navier--Stokes regularity problem, the fluid
obeys
\[
  \partial_tu+(u\cdot\nabla)u+\nabla p=\nu\Delta u,
  \qquad
  \nabla\cdot u=0,
  \qquad
  u(\cdot,0)=u_0.
\]
A strain in the flow pulls a vortex tube along its axis. As the vortex
stretches, incompressibility forces it to contract across its width. Each
material parcel keeps its volume, so lengthening the tube narrows its
cross-section and the rotating core inside it. That strain also acts on
vorticity \(\omega=\nabla\times u\). Taking the curl of the equation gives
\[
  \partial_t\omega+(u\cdot\nabla)\omega
  =
  (\omega\cdot\nabla)u+\nu\Delta\omega.
\]
The left side follows the moving fluid. When vorticity and strain align, the
stretching term is positive: \(|\omega|\) rises while the core narrows. Since
\(\omega\) is one spatial derivative of \(u\), the velocity changes more
sharply across that core. The gradients rise.

That narrowing brings the nonlinear and viscous parts of the equation into
direct competition. Transport can keep pulling the profile into a steeper
shape; viscosity spreads sharp differences through \(\nu\Delta u\). Scaling
shows what happens to this balance when the whole profile is compressed. For
\(\lambda>1\), set
\[
  u_\lambda(x,t)=\lambda u(\lambda x,\lambda^2t),
  \qquad
  p_\lambda(x,t)=\lambda^2p(\lambda x,\lambda^2t).
\]
Space has narrowed by \(1/\lambda\), while the corresponding time interval has
narrowed by \(1/\lambda^2\). Every term in the equation acquires a factor
\(\lambda^3\):
\[
\begin{aligned}
  \partial_tu_\lambda
  &=\lambda^3(\partial_tu)(\lambda x,\lambda^2t),\\
  (u_\lambda\cdot\nabla)u_\lambda
  &=\lambda^3((u\cdot\nabla)u)(\lambda x,\lambda^2t),\\
  \nabla p_\lambda
  &=\lambda^3(\nabla p)(\lambda x,\lambda^2t),\\
  \nu\Delta u_\lambda
  &=\lambda^3(\nu\Delta u)(\lambda x,\lambda^2t).
\end{aligned}
\]
Compression does not give either nonlinear transport or viscosity an extra
power of scale. It reproduces their balance inside a narrower region and a
shorter interval.

The Sobolev scale records that concentration by
\[
  \|u_\lambda(\cdot,t)\|_{\dot H^s}
  =
  \lambda^{s-\frac12}
  \|u(\cdot,\lambda^2t)\|_{\dot H^s}.
\]
Below \(s=\tfrac12\), the norm falls as the profile concentrates. Above
\(s=\tfrac12\), it rises. At \(s=\tfrac12\), it is unchanged. This is the
critical height: narrowing gains or loses no power of scale there. Kinetic
energy sits below it. A smaller region can
carry steeper gradients while its \(L^2\) cost remains finite, so finite energy
alone does not close the concentration route.

The time scaling makes the possible cascade just as sharp. If the \(k\)-th
stage has spatial width \(\lambda^{-k}\), its natural duration is
\(\lambda^{-2k}\). Those durations can accumulate:
\[
  \sum_{k=0}^{\infty}\lambda^{-2k}
  =
  \frac{1}{1-\lambda^{-2}}
  <
  \infty.
\]
An infinite sequence of narrower, faster stages could then fit below one
finite time. This is the Zeno cascade the proof has to rule out.

A finite maximal time can persist only if a continuation norm fails; for the
smooth history used here, bounded \(H^3\) control of \(u\) would already
continue the solution. If \(\partial_t^{j+1}u\) stayed integrable up to that
time in a fixed Sobolev norm, temporal integration would give
\(\partial_t^ju\) a finite endpoint. Read backward, each failing response
forces the next one out of that integrable control.

In order for a singularity to form, velocity has to blow up.

In order for that to happen, its acceleration would also have to blow up,
which means its jerk has to blow up, and so on, and so on:
\[
  u,\qquad
  \partial_tu,\qquad
  \partial_t^2u,\qquad
  \partial_t^3u,\qquad\ldots
\]
Following that sequence upward, one would expect less spatial regularity, not
more. A response changing faster in time should look rougher in space. The
equation reveals the reverse relation.

The first response is
\[
  \partial_tu
  =
  \nu\Delta u
  -(u\cdot\nabla)u
  -\nabla p.
\]
The next two are
\[
\begin{aligned}
  \partial_t^2u
  &=
  \nu\Delta\partial_tu
  -(\partial_tu\cdot\nabla)u
  -(u\cdot\nabla)\partial_tu
  -\nabla\partial_tp,
\end{aligned}
\]
and
\[
\begin{aligned}
  \partial_t^3u
  &=
  \nu\Delta\partial_t^2u
  -(\partial_t^2u\cdot\nabla)u\\
  &\quad
  -2(\partial_tu\cdot\nabla)\partial_tu
  -(u\cdot\nabla)\partial_t^2u
  -\nabla\partial_t^2p.
\end{aligned}
\]
The product rule has produced one transport term, then two, then three, with
the middle term counted twice. At every finite order the complete binomial
pattern is
\[
\begin{aligned}
  \partial_t^{r+1}u
  &=
  \nu\Delta\partial_t^ru\\
  &\quad
  -\sum_{a=0}^{r}\binom ra
  (\partial_t^au\cdot\nabla)\partial_t^{r-a}u
  -\nabla\partial_t^rp.
\end{aligned}
\]
Only now is there a pattern worth compressing. Write
\[
  V_r:=\partial_t^ru,
  \qquad
  Q_r:=\partial_t^rp.
\]
Then
\[
  V_{r+1}
  =
  \nu\Delta V_r
  -\sum_{a=0}^{r}\binom ra
  (V_a\cdot\nabla)V_{r-a}
  -\nabla Q_r,
\]
while incompressibility fixes the pressure coordinate through
\[
  Q_r
  =
  R_iR_j\sum_{a=0}^{r}\binom ra V_{a,i}V_{r-a,j}.
\]
Velocity and pressure travel through the ladder together. In each new temporal
response, \(\nu\Delta V_r\) reaches two spatial derivatives into the preceding
response. Transport and pressure carry the coupled terms, but the reversal is
already exposed: rising temporal order is tied to rising spatial order.

The critical energy shows exactly where those higher spatial orders sit. Let
\(\Lambda=(-\Delta)^{1/2}\) and
\[
  E_s(t):=\frac12\|\Lambda^su(t)\|_2^2,
  \qquad
  H(t):=E_{1/2}(t).
\]
Pairing the equation at height \(s\) gives
\[
  E_s'=\mathcal P_s-2\nu E_{s+1},
\]
where \(\mathcal P_s\) is the complete nonlinear transfer at that height. At
the critical height,
\[
\begin{aligned}
  H'
  &=\mathcal P_{1/2}-2\nu E_{3/2},\\
  H''
  &=\partial_t\mathcal P_{1/2}
    -2\nu\mathcal P_{3/2}
    +(2\nu)^2E_{5/2},\\
  H'''
  &=\partial_t^2\mathcal P_{1/2}
    -2\nu\partial_t\mathcal P_{3/2}
    +(2\nu)^2\mathcal P_{5/2}
    -(2\nu)^3E_{7/2}.
\end{aligned}
\]
Each step upward in time exposes the next Sobolev height. At order \(n\),
\[
  H^{(n)}
  =
  (-2\nu)^nE_{n+\frac12}
  {}+
  \sum_{j=0}^{n-1}
  (-2\nu)^j
  \partial_t^{\,n-1-j}\mathcal P_{j+\frac12}.
\]
\newpage
Here is the reversal. The imagined singular cascade predicts less spatial
regularity as its temporal responses accelerate. Viscosity links those
responses in the opposite direction: the \(n\)-th temporal level exposes the
spatial level \(n+\tfrac12\). Carry every finite temporal level and every
finite interaction term to the terminal time, and every finite Sobolev height
is present:
\[
  \bigcap_{m<\infty}H^m(\R^3)
  =
  H^\infty(\R^3)
  \hookrightarrow
  C^\infty(\R^3).
\]
Velocity, acceleration, jerk, and all later responses have led straight back
to smoothness.

The identities reveal the coordinates the proof needs. The finite rectangle
theorem supplies their whole-terminal control. Fix any finite temporal depth
\(N\) and spatial height \(M\). For every \(0\leq n\leq N\), it controls the
adjacent pair
\[
  V_n\in L^2(0,T_*;H^{M+1}),
  \qquad
  V_{n+1}\in L^2(0,T_*;H^{M-1}),
\]
together with its pressure coordinates. Equivalently, the finite rectangle
\[
  \cR_{N,M}(u;T)
  :=
  \sum_{n=0}^{N}
  \left(
    \|V_n\|_{L^2(0,T;H^{M+1})}^2
    {}+
    \|V_{n+1}\|_{L^2(0,T;H^{M-1})}^2
  \right)
\]
satisfies
\[
  \sup_{0<T<T_*}\cR_{N,M}(u;T)<\infty.
\]
The bound may depend on \(N\) and \(M\); it does not depend on the cutoff
\(T<T_*\). Each rectangle is finite, pressure-complete, and drawn from one
fixed-viscosity maximal history.

A larger rectangle retains every coordinate shared with a smaller one. The
overlaps agree because they are restrictions of that one history. Their
compatible projective family is \(\I[u_0]\), and for every finite pair
\((r,m)\) it gives
\[
  u\in H^r(0,T_*;H^m(\R^3)).
\]
No infinite-order bound is being demanded. Given a finite coordinate, one
chooses a rectangle large enough to contain it. Compatibility puts all such
choices into one intersection.

Now take the zeroth temporal row, denoted by \(\Szero\). For each finite \(m\),
\[
  u\in L^2(0,T_*;H^{m+2}).
\]
The pressure-completed interaction and the viscous term then belong to
\(L^1(0,T_*;H^m)\), so the equation gives
\[
  u_t\in L^1(0,T_*;H^m),
  \qquad
  u\in W^{1,1}(0,T_*;H^m)
  \hookrightarrow C([0,T_*];H^m).
\]
The endpoint values at different heights agree because they are limits of one
velocity field. They form one terminal state
\[
  u_*:=u(T_*)\in\bigcap_{m<\infty}H_\sigma^m
  =
  H_\sigma^\infty
  \hookrightarrow C^\infty.
\]
At rank \(r\), the pressure formula gives \(Q_r^*\) from
\(V_0^*,\ldots,V_r^*\), and the recurrence then gives \(V_{r+1}^*\).
Induction produces the full pressure-complete terminal jet. Smooth local theory
restarts the equation from \(u_*\); uniqueness glues that restart onto the
original history. A finite maximal time would then have been extended, which
is impossible. Hence \(T_*=\infty\).

This is not an a priori estimate proof. It is an intersection-to-estimate
proof. Viscosity relates the velocity's time derivatives to its spatial
Sobolev derivatives. Taken through every finite order, those two directions
meet in a common intersection. At that intersection,
\(H^\infty(\R^3)\hookrightarrow C^\infty(\R^3)\) gives smoothness, and the
familiar estimates can be taken from whichever finite Sobolev height they
require. The estimates come from the intersection, not the other way around.

For a finite spatial derivative order \(k\) and \(2\leq p\leq\infty\), choose
\[
  m>k+\frac32-\frac3p.
\]
That finite coordinate is already present, and Sobolev embedding gives
\[
  \I[u_0]
  \longrightarrow
  C_tH_x^m
  \longrightarrow
  L_t^\infty W_x^{k,p}.
\]
Choosing \(m=1,2,3\) recovers, in succession,
\[
\begin{aligned}
  \I[u_0]&\longrightarrow C_tH_x^1
    \longrightarrow L_t^\infty L_x^6,\\
  \I[u_0]&\longrightarrow C_tH_x^2
    \longrightarrow L_t^\infty L_x^\infty,\\
  \I[u_0]&\longrightarrow C_tH_x^3
    \longrightarrow L_t^\infty W_x^{1,\infty}.
\end{aligned}
\]
The proof has not climbed toward one privileged estimate. It has built the
intersection from which each finite demand can be read.

In the paper's compact notation, the whole motion is
\[
\begin{aligned}
  u_0
  &\longrightarrow
  \{(V_n,Q_n)\}_{n\geq0}
  \longrightarrow
  \{\cR_{N,M}[u_0;T_*]\}_{N,M<\infty}
  \longrightarrow
  \I[u_0]\\
  &\longrightarrow
  \Szero
  \longrightarrow
  u_*\in H^\infty
  \longrightarrow
  \{(V_n^*,Q_n^*)\}_{n\geq0}
  \longrightarrow
  T_*=\infty.
\end{aligned}
\]

\end{document}
